\par\SourceLeaf{138}
Bei den Universitätsseitig gemachten Ueberschlägen hat man jene 3 oder \(2\tfrac12\) Ehen gegen Eine Witwe für zu viel gehalten, und das Verhältniss von 2 Ehen gegen Eine Witwe zum Grunde gelegt.

Ich habe die Stelle des P. M. hier bloss deswegen angeführt, weil dadurch erklärlich wird, dass man sich bei dem Verhältniss von Zwei Ehen gegen Eine Witwe so leicht beruhigt hat, obgleich, sehr wahrscheinlich, auch dieses den Verhältnissen der Professoren-Witwenkasse noch nicht angemessen ist, sondern noch weniger Ehen gegen Eine Witwe gerechnet werden sollten. Ueber die Sache selbst wird das Nähere weiter unten vorkommen; aber der Geschäftsverlauf erinnert (wenn es erlaubt ist, ein Gleichniss aus einer niedern Sphäre hieher zu ziehen) unwillkürlich an Käufer, die bei unvollkommener eigner Waarenkenntniss einen guten Handel gemacht zu haben glauben, wenn sie weit unter dem zuerst geforderten Preise eingekauft haben, obgleich sie, bei Lichte besehen, noch immer zu theuer bezahlten. Ich brauche nicht zu erinnern, dass ich diess Gleichniss nicht über die Gebühr ausgedehnt wissen will, denn der unbekannte Proponent hat die Verhältnisse \(3:1\) und \(2\tfrac12:1\) ohne Zweifel in gutem Glauben an ihre Zulässigkeit vorgebracht.

Indem der damalige Prorector F., unter dem 6. Julius, das Rescript bei dem Senate in Umlauf setzt, fügt er den beiden darin enthaltenen Deliberationsgegenständen (Erhöhung der Beiträge und Erhöhung der Pensionen) noch einen dritten bei, durch den Vorschlag, die Dauer der Waisenpensionen bis zum vollendeten \(20^{\text{sten}}\) Lebensjahre zu erweitern. Er überlässt den Senatsmitgliedern, sich über diese Gegenstände gleich schriftlich, oder in der auf den 12. Julius angesetzten Senatsversammlung zu äussern.

Diese Missive ist von 17 Senatsmitgliedern unterzeichnet; von den dabei gefallenen Aeusserungen sind hier nur ein paar zu erwähnen.

Der damalige Curator der Witwenkasse, P., stellt vor allem den Grundsatz auf: die Kasse sei den gegenwärtigen Witwen eben so viel schuldig als den künftigen, sie sei aber auch den künftigen Witwen genau so viel schuldig wie den gegenwärtigen. Diese (an sich in der That sehr vage) Phrase erläutert er dahin, dass jede künftige Witwe, welche weniger erhalte, als eine andere früher erhalten habe, (seiner Meinung nach) wahrhaft lädirt werde, und das erste Princip müsse demnach sein, die Pensionshöhe so zu bestimmen, dass, nach höchster Wahrscheinlichkeit, sie niemals wieder vermindert zu werden brauche. In dieser Beziehung hält aber P. den Calcül in der Beilage des Rescripts nicht für sicher genug; man dürfe nicht \(2\tfrac12\) Ehen auf 1 Witwe, sondern nur 2 rechnen, und müsse also das Maximum der Witwen nicht auf 10—12 sondern auf 14—15 setzen, mithin auch geringere Pensionshöhen annehmen.

\textsc{Kaestner} hält die Frage für zu verwickelt und schwierig, als dass sich ohne eine \textit{umständliche} und \textit{genaue} Untersuchung etwas festsetzen lasse; auch er sei der Meinung, dass mehr nicht als \textit{höchstens} 2 Ehen auf 1 Witwe, gerechnet werden dürfen. Da er längst aus der Witwenkasse ausgeschieden sei (er war Theilnehmer gewesen von 1755—1773), so habe er in der Sache keine Stimme (als Senior der philosophischen Facultät war er doch Mitglied der Kirchendeputation), rathe aber, keinen Beschluss zu fassen, ohne vorher einen Sachverständigen, etwa den p. \textsc{Kritter} zu befragen.

G. wünscht auch, dass durch genaue Rechnungen die Kräfte der Kasse ermittelt werden möchten, und weiset auf die \textit{schrecklichen} Folgen übereilter Beschliessung zu grosser Pensionen, an den Beispielen der Hannoverschen, Bremischen u. a. Witwengesellschaften hin.

Die übrigen Vota stimmen theils den vorigen bei, theils entwickeln sie Bedenklichkeiten, wegen Erhöhung der Beiträge oder Verlängerung der Waisenpensionen, was hier nicht extrahirt zu werden braucht.

In der Senatssitzung vom 12. Julius, in welcher, den Prorector mitgezählt, 12 Professoren anwesend waren, erklärten sich \textit{für} die Erhöhung der Beiträge 9 unbedingt, 2 mit dem Zusatz, dass sie die Erhöhung für zu gross hielten; einer (der bei der schriftlichen Votirung sich nachdrücklich dagegen erklärt hatte) wollte den mehrsten Stimmen beitreten. — Der F.’sche Vorschlag, wegen Verlängerung der Waisenpensionen wurde bis zu genauerer Erwägung der Umstände der Kasse noch beanstandet. — Wegen
